\documentclass[../../../hsm_tok_q.tex]{subfiles}

\begin{document}
    \setcounter{figure}{0}
    \textsf{図\ref{fig/hsm_tok_2018_2nd_04_01_q}}で，%
    四角形ABCDは，平行四辺形である．

    点Pは，辺BC上にある点で，頂点B，頂点Cのいずれにも一致しない．

    点Aと点Pを結ぶ．
    
    次の各問に答えよ    \\
    \begin{minipage}{\linewidth}
        \vspace{.5\baselineskip}
        {\centering
        \setlength{\abovecaptionskip}{2pt}
        \includegraphics%
            {\subfix{fig_hsm_tok_2018_2nd_04/fig_hsm_tok_2018_2nd_04_01_q.pdf}}
            \captionof{figure}{} \label{fig/hsm_tok_2018_2nd_04_01_q}
        }
    \end{minipage}
    \vspace{.1\baselineskip}
    %
    \begin{enumerate}
        \item \textsf{図\ref{fig/hsm_tok_2018_2nd_04_01_q}}において，%
            $\angle\mathrm{ABC}$＝$50^\circ$，%
            $\angle\mathrm{BAP}$＝$a^\circ$とするとき，%
            $\angle\mathrm{DAP}$の大きさを表す式を，%
            次の\textsf{ア}～\textsf{エ}のうちから選び，記号で答えよ．

            \vspace{.15\baselineskip}
            \begin{tabularx}{\dimexpr\linewidth-3.5mm}{@{}LLLL@{}}
                {\textsf ア}\quad$(50-a)$度 & {\textsf イ}\quad$(a+50)$度 \\ [1pt]
                {\textsf ウ}\quad$(130-a)$度 & {\textsf エ}\quad$(a+130)$度
            \end{tabularx}

        \item \textsf{図\ref{fig/hsm_tok_2018_2nd_04_02_q}}は，%
            \textsf{図\ref{fig/hsm_tok_2018_2nd_04_01_q}}において，%
            線分APと線分BDとの交点をQとした場合を表している．

            次の①，②に答えよ．
            %
            \begin{enumerate}
                \item[①] $\mathrm{△AQD}$∽$\mathrm{△PQB}$であることを証明せよ．

                \item[②] 次の\:\myboxa{　　　}\:の中の「\textsf{う}」「\textsf{え}」%
                    「\textsf{お}」に当てはまる数字をそれぞれ答えよ．

                    \textsf{図\ref{fig/hsm_tok_2018_2nd_04_02_q}}において，%
                    点Pを通り線分BDと平行な直線を引き，辺CDとの交点をR，%
                    頂点Aと点Rを結び，線分ARと線分BDとの交点をSとし，%
                    点Pと点Sを結んだ場合を考える．

                    BP：PC＝3：2のとき，$\mathrm{△PSQ}$の面積は，%
                    四角形ABCDの面積の$\dfrac{\,\myboxa{　\textsf{う}　}\,}%
                    {\myboxa{\hspace{.45\zw}\textsf{え\hspace{.1\zw}お}\hspace{.45\zw}}}$倍である．
            \end{enumerate}
            %
            \begin{minipage}{\linewidth}
                \vspace{.5\baselineskip}
                \par{\centering
                    \setlength{\abovecaptionskip}{3pt}
                    \includegraphics%
                        {\subfix{fig_hsm_tok_2018_2nd_04/fig_hsm_tok_2018_2nd_04_02_q.pdf}}
                        \captionof{figure}{} \label{fig/hsm_tok_2018_2nd_04_02_q}
                \par}
            \end{minipage}            
        \end{enumerate}

\end{document}


